\pdfoutput=1
\documentclass[11pt]{amsart}
\usepackage{amsmath,amssymb,amsthm,booktabs}
\usepackage[margin=1.15in]{geometry}
\usepackage[colorlinks=true,linkcolor=blue,citecolor=blue,urlcolor=blue]{hyperref}
\hypersetup{pdfauthor={Bernd Johannes Wuebben},
pdftitle={Small Dolbeault eigenvalues near stability walls}}
\newtheorem{theorem}{Theorem}[section]
\newtheorem{proposition}[theorem]{Proposition}
\newtheorem{lemma}[theorem]{Lemma}
\newtheorem{corollary}[theorem]{Corollary}
\theoremstyle{remark}
\newtheorem{remark}[theorem]{Remark}
\newtheorem{example}[theorem]{Example}
\newcommand{\End}{\operatorname{End}}
\newcommand{\Hom}{\operatorname{Hom}}
\newcommand{\Ext}{\operatorname{Ext}}
\newcommand{\tr}{\operatorname{tr}}
\newcommand{\rk}{\operatorname{rk}}
\newcommand{\id}{\operatorname{id}}
\newcommand{\diag}{\operatorname{diag}}
\newcommand{\PP}{\mathbb P}
\newcommand{\CC}{\mathbb C}
\newcommand{\RR}{\mathbb R}
\newcommand{\HH}{\mathcal H}
\newcommand{\db}{\bar\partial}
\title{Small Dolbeault eigenvalues near stability walls}
\author{Bernd Johannes Wuebben}
\subjclass[2020]{53C07 (primary); 53C55, 58J50 (secondary)}
\keywords{Hermitian--Einstein metric, slope stability, Dolbeault Laplacian,
wall crossing, vector bundle extension, toric variety}
\date{October 27, 2026}

\begin{document}
\begin{abstract}
We determine the small spectrum of the Dolbeault Laplacian on trace-free
endomorphisms of Hermitian--Einstein bundles near intersections of
slope-stability walls. On a compact K\"ahler manifold of dimension at
least two, consider an extension of $\bigoplus_{j=1}^kF_j$ by $F_0$,
where the factors are pairwise nonisomorphic and stable of equal slope
at the limiting K\"ahler class, and every extension component is nonzero.
Write $r_i=\rk F_i$ and let $\epsilon_j$ be the Einstein constant of
$F_j$ minus that of the extension. In the chamber $\epsilon_j>0$,
exactly $k$ eigenvalues tend to zero. They agree, to relative error
$O(\sum_j\epsilon_j)$, with the eigenvalues of
$\diag(\epsilon_1,\ldots,\epsilon_k)+r_0^{-1}vv^*$,
where $v_j=\sqrt{r_j\epsilon_j}$.
The estimate is uniform in all positive ratios of the slope gaps.
Thus the ranks and slopes determine the leading spectrum even when
the gaps vanish at arbitrarily different rates.

For two factors we obtain
$\lambda_1=\delta-D\delta^2+o(\delta^2)$, where $\delta>0$ is the
quotient Einstein constant minus the subbundle Einstein constant and
$D\ge0$ is given by the Green operator of the graded connection.
After rescaling the extension and making unitary identifications with
the graded limit, we determine the normalized first eigensection through
order $\delta$ in every spatial $C^m$ norm. Toric tangent bundles provide
examples with arbitrarily many small eigenvalues whose orders of
vanishing can be prescribed independently.
\end{abstract}
\maketitle

\section{Introduction}

When a stable bundle approaches a strictly semistable polarization,
its Hermitian--Einstein connection can converge, after gauge, to the
connection on its graded object. The additional endomorphisms of the
graded bundle then appear as small eigenvalues of the endomorphism
Dolbeault Laplacian. We ask how those eigenvalues depend on the distances
to the stability walls, especially when several distances tend to zero
at different rates.

Let $X$ be a compact connected K\"ahler manifold of complex dimension
$n\ge2$. For a K\"ahler form $\omega$ we use
\[
 \deg_\omega F=\int_Xc_1(F)\frac{\omega^{n-1}}{(n-1)!},\qquad
 \mu_\omega(F)=\frac{\deg_\omega F}{\rk F},\qquad
 V_\omega=\int_X\frac{\omega^n}{n!}.
\]
The Einstein constant is $\kappa_\omega(F)=2\pi\mu_\omega(F)/V_\omega$.
Our operator is $\Delta=\db_{\End}^*\db_{\End}$ on complex
trace-free endomorphisms. Eigenvalues are counted with complex
multiplicity. A scalar multiple of a bundle metric induces the same
endomorphism metric, so the spectrum of a stable bundle with its
Hermitian--Einstein metric is well defined once $\omega$ is fixed.

Fix an extension of holomorphic vector bundles
\begin{equation}\label{eq:star}
 0\longrightarrow F_0\longrightarrow E
 \longrightarrow\bigoplus_{j=1}^kF_j\longrightarrow0.
\end{equation}
Assume that the $F_i$ are pairwise nonisomorphic, stable of the same
slope at $\omega_0$, and that every component of the extension class
in $\Ext^1(F_j,F_0)$ is nonzero. Write $r_i=\rk F_i$ and
$r=\sum_{i=0}^kr_i$. The extension graph has one central vertex and
$k$ leaves; we call \eqref{eq:star} a star extension.

Let $\omega_s$ be a smooth family of K\"ahler forms, with $s$ in a
neighborhood of zero in $\RR^d$. Put
\begin{equation}\label{eq:epsilon}
 \epsilon_j(s)=\kappa_{\omega_s}(F_j)-\kappa_{\omega_s}(E),
 \qquad e(s)=\sum_{j=1}^k\epsilon_j(s).
\end{equation}
These are slope surpluses of the quotient factors. Stability requires
$\epsilon_j>0$: the kernel of $E\to F_j$ must have smaller slope
than $E$.

\begin{theorem}\label{thm:main}
Under the preceding assumptions, for $s$ sufficiently close to zero,
$E$ is stable with respect to $\omega_s$ if and only if
$\epsilon_j(s)>0$ for every $j$.
In this region let $\lambda_1(s)\le\lambda_2(s)\le\cdots$ be the
endomorphism Dolbeault eigenvalues of its Hermitian--Einstein metric.
Let $\nu_1(s)\le\cdots\le\nu_k(s)$ be the eigenvalues of
\begin{equation}\label{eq:rank-matrix}
 S(\epsilon)=\diag(\epsilon_1,\ldots,\epsilon_k)
       +\frac1{r_0}vv^*,\qquad v_j=\sqrt{r_j\epsilon_j}.
\end{equation}
There are constants $C,\gamma>0$ such that
\begin{equation}\label{eq:main-relative}
 \left|\frac{\lambda_j(s)}{\nu_j(s)}-1\right|\le Ce(s)
 \quad(1\le j\le k),\qquad \lambda_{k+1}(s)\ge\gamma.
\end{equation}
The constants are independent of the positive ratios
$\epsilon_i/\epsilon_j$. No transversality of the parameter family
to the walls is required.
\end{theorem}

The uniformity in \eqref{eq:main-relative} is a relative assertion for
each small eigenvalue. An additive error controlled by the largest gap
would not determine an eigenvalue associated with a much smaller gap.
Here one may, for example, let one gap be exponentially small in the
reciprocal of another. The constants may depend on the fixed bundles,
extension, parameter family and number of factors.

The matrix in \eqref{eq:rank-matrix} is the weighted star-graph
Laplacian written on the complement of its constant mode. Its Rayleigh
quotient is
\begin{equation}\label{eq:graph-intro}
 \frac{\sum_{j=1}^kr_j\epsilon_j|x_j-x_0|^2}
      {\sum_{i=0}^kr_i|x_i|^2},\qquad \sum_{i=0}^kr_ix_i=0.
\end{equation}
When all surpluses equal $\epsilon$, its eigenvalues are $\epsilon$
with multiplicity $k-1$ and $(r/r_0)\epsilon$ once. The harmonic
norms of the extension classes do not occur in this expression.

The proof first solves the Hermitian--Einstein equation orthogonally
to the block-scalar endomorphisms. If $t_j$ scales the $j$th harmonic
extension and $w_j=t_j^2$, the remaining equations have the form
\[
 V_sr_j\epsilon_j-w_jb_j(s,w)=0,
 \qquad b_j(s,0)=\|\beta_{j,s}\|^2>0.
\]
The factor $w_j$ in its own equation is essential. On the face $t_j=0$
the $j$th summand splits off, and its integrated curvature is fixed by
its degree. Equivariance and smoothness therefore give exact
divisibility by $w_j$, which yields a relative estimate for every
extension weight without dividing by the smallest slope gap. A relative
Schur estimate identifies the harmonic-family spectrum with the graph
spectrum. A pointwise comparison of the corrected and split metrics
then transfers the estimate to the actual Hermitian--Einstein metric.

For two factors write
$\delta=\kappa_{\omega_s}(F_1)-\kappa_{\omega_s}(F_0)$.
Theorem~\ref{thm:second-order} refines the leading result to
\[
 \lambda_1=\delta-D\delta^2+o(\delta^2),\qquad D\ge0,
\]
and gives a smooth eigensection expansion. The subbundle projection
has Rayleigh quotient exactly $\delta$ at an HE metric. This identity
fixes the first term without an unknown metric correction and makes
the second term a Green-operator expression at the graded limit.

\subsection*{Relation to previous work}
The Hermitian--Einstein correspondence supplies the metrics on the
stable factors and identifies simple HE bundles with stable bundles
\cite{UY,Kobayashi}. Clarke--Tipler \cite{CT} prove local chamber
decomposition and convergence of HYM connections to the graded
connection for simple bundles with locally free graded object and
abelian graded automorphism group. Their Proposition~3.2 gives a
parameter-dependent complementary implicit-function construction;
their Theorem~1.2 gives convergence along paths in a stable chamber.
The bundles considered here satisfy those hypotheses. Delloque
\cite{Delloque} extends this continuity theory to balanced metrics
without the abelian automorphism hypothesis. His Theorem~4 also gives
a norm bound for the difference from the graded connection.
We use a specialized construction to determine the small spectrum
quantitatively.

Sektnan--Tipler \cite{ST} construct HYM connections on pullback bundles
in adiabatic classes. Section~5.2 develops scaled
harmonic extensions and balances their curvature against slope defects.
Proposition~50 (Proposition~5.27 of the preprint) supplies a lower bound
for the linearized operator;
these are substantial precedents for the construction used here.
Our estimates concern all small eigenvalues near an interior K\"ahler
class and are relative uniformly in the independent positive wall
distances.

The earlier spectral paper \cite{Flow} gives relative graph comparison
for harmonic families and quantitative transfer along HYM flow. Its
two-factor harmonic family has expansion $ct^2-dt^4+O(t^6)$.
The coefficient in our HE expansion is exactly $D=d/c^2$ after
normalization. The new assertion concerns its occurrence for the
actual equilibrium metrics, with the slope gap as parameter.
For several factors, the principal result is the cancellation of
extension norms from the leading HE spectrum, together with the
uniform relative error across arbitrarily separated wall distances.
The proof does not take a joint limit in polarization and flow time.

The toric companion \cite{Chambers} supplies tangent-bundle extensions
of the form \eqref{eq:star}, including examples with arbitrarily many
independent slope coordinates. Section~\ref{sec:toric} states the
geometric inputs and applies the analytic theorem. The analytic proof
does not use a finite toric census.

\section{The Hermitian--Einstein equation near the graded bundle}
\label{sec:construction}

We retain the hypotheses of Theorem~\ref{thm:main}. Norm estimates in
this section are uniform after shrinking the parameter neighborhood.
They do not compare $|s|$ with any positive slope surplus.

\subsection{Factor metrics and harmonic extensions}
Choose smoothly varying HE metrics $h_{i,s}$ on the fixed holomorphic
bundles $F_i$, with scalar normalizations. Here is the local analytic
justification. First choose normalized HE determinant metrics, using
the scalar Laplace equation. With those determinants fixed, the
linearization of the trace-free HE equation at $s=0$ is the
endomorphism Dolbeault Laplacian on trace-free Hermitian sections.
It is invertible because $F_i$ is stable. For an integer $\ell>n$,
the equation is a smooth map from $H^{\ell+2}$ to $H^\ell$.
The implicit-function theorem and elliptic regularity give the
metrics for small $s$. Their existence and simplicity also show that
the factors remain stable.

Let $\beta_{j,s}$ be the harmonic representative of the fixed nonzero
class in $H^{0,1}(\Hom(F_j,F_0))$. These representatives vary
smoothly. Indeed the holomorphic Hom bundle is fixed, so the
cohomology dimension is constant; the spectral projection onto the
kernel of its smoothly varying elliptic Laplacian is smooth. Elliptic
regularity gives smoothness in every spatial norm. Write
\[
 B_j(s)=\|\beta_{j,s}\|_s^2>0,\qquad
 h_{\mathrm{split},s}=\bigoplus_{i=0}^kh_{i,s}.
\]
Each $B_j$ has a uniform positive lower bound.

On the fixed smooth direct sum define
\begin{equation}\label{eq:scaled-dbar}
 \db_{s,t}=\db_0+\sum_{j=1}^kt_j\beta_{j,s},
 \qquad t\in\RR^k,
\end{equation}
where $\beta_{j,s}$ occupies the block $F_j\to F_0$.
The operator is integrable in every complex dimension:
$\db_0\beta_{j,s}=0$, and every product of two off-diagonal
entries vanishes. When all $t_j\ne0$, independent constant
scalings of the quotient factors identify this extension with
\eqref{eq:star}. Changes of harmonic representative are implemented
by upper-triangular complex gauges. Its determinant holomorphic
structure is independent of $t$.

\subsection{The complementary equation}
Let $\HH$ be the real space of constant block-scalar Hermitian
endomorphisms with total trace zero. It has dimension $k$. Let
$P_s$ be its $L^2_s$ projection and $Q_s=I-P_s$, acting on
pointwise trace-free Hermitian sections. We seek a metric
$h=h_{\mathrm{split},s}e^u$, where $u$ is a trace-free Hermitian
endomorphism satisfying $P_su=0$.
Express its mean curvature in the fixed Hermitian structure:
\begin{equation}\label{eq:Khat}
 \widehat K(s,t,u)
 =e^{u/2}\bigl(i\Lambda_{\omega_s}F_{\db_{s,t},h_{\mathrm{split},s}e^u}\bigr)e^{-u/2}.
\end{equation}
Because $\tr u=0$, its trace is the constant $r\kappa_{\omega_s}(E)$.

\begin{lemma}\label{lem:complement}
For small $(s,t)$ there is a unique small normalized $u(s,t)$ solving
\[
 Q_s\bigl(\widehat K(s,t,u)-\kappa_{\omega_s}(E)I\bigr)=0.
\]
It is smooth in the parameters and spatial variables, and for every $m$
\begin{equation}\label{eq:u-estimate}
 \|u(s,t)\|_{C^m}\le C_m|t|^2.
\end{equation}
The construction is equivariant under independent constant unitary
scalings of the factors.
\end{lemma}

\begin{proof}
Use the real Hermitian spaces $H^{\ell+2}\to H^\ell$, $\ell>n$.
Block-diagonal positive changes of metric identify the Hermitian bundles
for different $s$, preserve trace and block scalars, and are bounded
with bounded inverses in every fixed Sobolev norm. The projections
$Q_s$ depend smoothly on $s$. After those identifications,
$Q_s|_{\HH^{\perp_0}}$ is an isomorphism onto
$\HH^{\perp_s}$ for small $s$. These maps give fixed Banach spaces
for the parameter-dependent equation.

At $t=u=0$, the split curvature $K_{0,s}$ is block scalar, with
entries $\kappa_i(s)I_{F_i}$. Differentiation of the un-conjugated
mean curvature in the convention $h=h_{\mathrm{split},s}e^u$ gives
$\Delta_{\partial,s}u$. Differentiating the conjugation in
\eqref{eq:Khat} adds $[u,K_{0,s}]/2$. The K\"ahler identity on
endomorphism-valued zero-forms is
$\Delta_{\partial}-\Delta_{\db}=[K_{0,s},\cdot]$. Thus the real
Hermitian linearization is
\begin{equation}\label{eq:real-linearization}
 L_s=\Delta_{\db,s}+\tfrac12[K_{0,s},\cdot]
    =\tfrac12(\Delta_{\db,s}+\Delta_{\partial,s}).
\end{equation}
It is self-adjoint and nonnegative and annihilates $\HH$.
At $s=0$ all factor Einstein constants agree, so
$L_0=\Delta_{\db,0}$. Its trace-free Hermitian kernel is exactly
$\HH$: each factor has only scalar endomorphisms, and the Hom groups
between distinct stable factors of equal slope vanish.

Consequently $Q_0L_0$ is an isomorphism from the complementary
$H^{\ell+2}$ space to $H^\ell$. In fixed spaces, $Q_sL_s$ converges
in this operator norm to $Q_0L_0$. The Neumann-series criterion gives
uniformly bounded inverses. The implicit-function theorem now produces
$u(s,t)$, and uniqueness gives $u(s,0)=0$.

The derivative of the split mean curvature in $t_j$ has upper
off-diagonal entry $-\db_{\Hom}^*\beta_{j,s}$ and lower entry
its Hermitian adjoint. Both vanish by harmonicity for every $s$.
Differentiating the complementary equation therefore
gives $u_{t_j}(s,0)=0$, proving \eqref{eq:u-estimate}. Applying the
IFT at higher Sobolev indices and using uniqueness, or bootstrapping
the equation, gives each $C^m$ estimate. All equations, projections
and normalizations are invariant under the indicated unitary scalings;
uniqueness proves equivariance.
\end{proof}

\subsection{Scalar equations on the coordinate faces}
Set $p_j=I_{F_j}-(r_j/r)I$ for $1\le j\le k$. These form a basis
of $\HH$. The remaining curvature equations are
\[
 f_j(s,t)=\bigl\langle p_j,
 \widehat K(s,t,u(s,t))-\kappa_{\omega_s}(E)I\bigr\rangle_s=0.
\]
Changing the sign of any $t_i$ is induced by a unitary sign on $F_i$,
and fixes every $p_j$. Thus every $f_j$ is separately even in the
$t_i$.

We use the following elementary smoothness fact, jointly with the
parameters. A smooth function $f(s,t)$ that is separately even in
$t_1,\ldots,t_k$ can be written $F(s,t_1^2,\ldots,t_k^2)$ with
$F$ smooth on the closed positive orthant and smoothly extendible
across its faces. In one variable, Taylor expansion has only even
powers, and its differentiated remainders give derivatives of every
order in $w=t^2$ up to $w=0$. A smooth extension from the half-line
then gives the assertion; apply this successively in each variable,
retaining the other variables and $s$ as parameters. If such an $F$
vanishes on $w_j=0$, the formula
\[
 F(s,w)=w_j\int_0^1\partial_{w_j}F(s,w_1,\ldots,\theta w_j,
                                      \ldots,w_k)\,d\theta
\]
gives smooth divisibility on the orthant. We write $w_j=t_j^2$ below.

\begin{proposition}\label{prop:face}
The scalar equations have the form
\begin{equation}\label{eq:scalar-factor}
 f_j(s,w)=V_sr_j\epsilon_j(s)-w_jb_j(s,w),\qquad
 b_j(s,w)=B_j(s)+O\left(\sum_iw_i\right).
\end{equation}
The functions $b_j$ are smooth jointly in the parameters and orthant
variables, and the estimates are uniform near the origin.
\end{proposition}

\begin{proof}
On the face $t_j=0$, $F_j$ splits holomorphically from the rest.
Its circle of scalar unitary transformations fixes the Dolbeault
operator, the background metric and the complementary equation.
Lemma~\ref{lem:complement} implies that $u(s,t)$ commutes with
this circle. The corrected metric therefore also splits off $F_j$.
Conjugation by $e^{u/2}$ preserves the trace of that block. Since
the total residual is pointwise trace-free, Chern--Weil gives exactly
\begin{equation}\label{eq:face-identity}
 f_j(s,t)\big|_{t_j=0}
 =2\pi\deg_{\omega_s}(F_j)-V_sr_j\kappa_{\omega_s}(E)
 =V_sr_j\epsilon_j(s).
\end{equation}
This equality is independent of all other extension parameters.
The smoothness fact just stated gives the factorization in
\eqref{eq:scalar-factor}.

It remains to determine $b_j(s,0)$. At the split metric the quadratic
curvature term of the entry $t_j\beta_{j,s}$ is positive on $F_0$
and negative on $F_j$. Its pairing with $p_j$ is
$-t_j^2\|\beta_{j,s}\|^2$. The diagonal terms from other edges pair
to zero, as do their off-diagonal cross terms. The quadratic metric
correction contributes $L_su_2$, whose pairing with $p_j$ vanishes
because $L_s$ is self-adjoint and $L_sp_j=0$. Hence
$b_j(s,0)=B_j(s)$. Smoothness yields the remainder uniformly in $s$.
\end{proof}

\begin{proposition}\label{prop:balance}
For small $s$ with every $\epsilon_j(s)>0$, there is a positive
solution $w(s)$ of the scalar equations, and
\begin{equation}\label{eq:individual-balance}
 w_j(s)B_j(s)=V_sr_j\epsilon_j(s)\bigl(1+O(e(s))\bigr)
 \qquad(1\le j\le k).
\end{equation}
The estimates are uniform in the positive gap ratios. The corresponding
corrected metric is an HE metric on $E$, and $E$ is stable.
\end{proposition}

\begin{proof}
Extend the $b_j$ smoothly to a full neighborhood of $w=0$, preserving
positivity after shrinking it. Introduce independent real targets
$\eta_j$ and solve
\[
 w_jb_j(s,w)=V_sr_j\eta_j.
\]
The $w$-Jacobian at $w=\eta=0$ is $\diag(B_j(s))$, uniformly
invertible. The parameter IFT gives $w(s,\eta)=O(|\eta|)$,
with $w(s,0)=0$. For positive targets the equations themselves force
each $w_j>0$, so the solution lies where the extension agrees with
the original functions. Set $\eta_j=\epsilon_j(s)$. Then
\[
 \frac{w_jB_j(s)}{V_sr_j\epsilon_j(s)}
       =\frac{B_j(s)}{b_j(s,w)}=1+O(e(s)).
\]
No estimate divides an error by $\min_j\epsilon_j$.

Together with the complementary equation, vanishing of the scalar
pairings against a basis of $\HH$ gives the full HE equation. Take
$t_j=\sqrt{w_j}$; all are nonzero, so the holomorphic bundle is $E$.
To see simplicity, an endomorphism of $E$ preserves $F_0$ because
$\Hom(F_0,F_j)=0$. Its actions on $F_0$ and on each quotient
$F_j$ are scalar. Compatibility with each nonzero extension class
forces all these scalars to agree. The possible remaining map from
the quotient to $F_0$ is zero because $\Hom(F_j,F_0)=0$.
Thus an HE metric makes $E$ stable. Necessity of the positive
surpluses follows from the kernels of $E\to F_j$, as in the introduction.
\end{proof}

\section{Relative spectral estimates}
\label{sec:spectrum}

We first use the metric $h_{\mathrm{split},s}$ on the holomorphic bundle
defined by \eqref{eq:scaled-dbar}. The comparison below is the
harmonic-family argument of \cite{Flow}; we give its proof to keep
track of uniformity in both the polarization and the extension weights.

\subsection{The harmonic family}
Write $D_0$ for the split Dolbeault operator on trace-free endomorphisms
and
\[
 D_t=D_0+A_t,\qquad A_tx=\left[\sum_jt_j\beta_{j,s},x\right].
\]
The complex space $\HH_\CC=\ker D_0$ consists of constant
block-scalar endomorphisms. Write them as
$x=\diag(x_iI_{F_i})$, where $\sum_ir_ix_i=0$.
Let $P$ be the $L^2_s$ projection onto this space and $Q=I-P$.
The spectral gap of $D_0^*D_0$ on its complement is uniformly
positive for small $s$. Since the harmonic forms vary smoothly,
$a_t=\|A_t\|_{L^2\to L^2}\le C|t|$.

For $x\in\HH_\CC$, harmonicity gives $D_0^*A_tx=0$. Distinct
matrix entries are orthogonal, so
\begin{equation}\label{eq:harmonic-graph}
 G_t(x):=\|D_tx\|^2
   =\sum_{j=1}^kw_jB_j(s)|x_j-x_0|^2,\qquad
 \|x\|^2=V_s\sum_{i=0}^kr_i|x_i|^2.
\end{equation}
No pointwise orthogonality among the forms themselves is required.
Let $g_1(t)\le\cdots\le g_k(t)$ be the eigenvalues of this
quadratic form with the displayed mass inner product.

\begin{lemma}\label{lem:relative-schur}
For small $(s,t)$ with all $t_j\ne0$, the first $k$ eigenvalues
of $\Delta_t=D_t^*D_t$ satisfy
\begin{equation}\label{eq:relative-schur}
 (1-C|t|^2)g_j(t)\le\lambda_j(h_{\mathrm{split},s},\db_{s,t})\le g_j(t).
\end{equation}
All remaining eigenvalues have a uniform positive lower bound.
The constant does not depend on any ratio of the nonzero $t_j$.
\end{lemma}

\begin{proof}
For $y\perp\HH_\CC$,
$\|D_ty\|\ge\|D_0y\|-a_t\|y\|$. Hence the compression
$C_t=Q\Delta_tQ$ is bounded below by a constant $\gamma>0$ for
small parameters. For $x\in\HH_\CC$ the harmonic cancellation
gives $\Delta_tx=A_t^*A_tx$. Its coupling to the complement,
$F_t=Q\Delta_tP$, therefore satisfies
\[
 \|F_tx\|^2\le a_t^2G_t(x).
\]
For $0\le\lambda<\gamma/2$ this implies the quadratic-form bound
\begin{equation}\label{eq:schur-form-bound}
 0\le F_t^*(C_t-\lambda)^{-1}F_t
       \le\frac{2a_t^2}{\gamma}G_t.
\end{equation}

The $k$ trial directions in $\HH_\CC$ give $k$ eigenvalues
$O(|t|^2)$ by min--max. The complementary subspace has codimension
$k$, so the next eigenvalue is at least $\gamma$. Below $\gamma/2$,
the Schur complement of $\Delta_t-\lambda$ is the finite operator
\[
 G_t-\lambda I-F_t^*(C_t-\lambda)^{-1}F_t
\]
on $\HH_\CC$. Because $C_t-\lambda$ is positive, the number of
negative eigenvalues of $\Delta_t-\lambda$ equals that of this
Schur complement. Comparison with
$(1-2a_t^2/\gamma)G_t-\lambda I$ and $G_t-\lambda I$ brackets
every ordered eigenvalue and proves \eqref{eq:relative-schur}.
The argument uses neither a lower bound on the smallest graph
eigenvalue nor a separation between graph eigenvalues.
\end{proof}

\subsection{The actual metric}
For the solution in Proposition~\ref{prop:balance},
Lemma~\ref{lem:complement} gives the pointwise metric comparison
\[
 e^{-C|w|}h_{\mathrm{split},s}\le h_{\mathrm{HE},s}
       \le e^{C|w|}h_{\mathrm{split},s}.
\]
If $e^{-a}h_0\le h_1\le e^ah_0$, the induced squared norms on
endomorphism-valued zero-forms and one-forms lie between $e^{-2a}$
and $e^{2a}$ times the original squared norms. The Dolbeault operator
and base K\"ahler form are the same on both sides. The Rayleigh
quotient, followed by min--max, therefore gives
\begin{equation}\label{eq:metric-comparison}
 e^{-4a}\lambda_j(h_0)\le\lambda_j(h_1)\le e^{4a}\lambda_j(h_0).
\end{equation}
These inequalities hold on the same algebraically trace-free section
space. They use no derivatives of the metric comparison.

Since $|w|=O(e(s))$, equations \eqref{eq:individual-balance},
\eqref{eq:relative-schur} and \eqref{eq:metric-comparison} show that
the first $k$ HE eigenvalues agree, to relative error $O(e(s))$,
with the graph form whose edge weights are $V_sr_j\epsilon_j$
and masses $V_sr_i$. Every weight comparison is relative, so it is
also a relative comparison of the entire graph quadratic form.
Min--max preserves that comparison for each ordered eigenvalue.
The common volume cancels, leaving \eqref{eq:graph-intro}.
The complementary spectral gap persists under the same metric
comparison.

To obtain \eqref{eq:rank-matrix}, set $y_j=x_j-x_0$. Eliminating
$x_0$ by the trace constraint gives mass matrix and energy matrix
\[
 M=\diag(r_1,\ldots,r_k)-\frac1r(r_ir_j)_{i,j=1}^k,
 \qquad D=\diag(r_1\epsilon_1,\ldots,r_k\epsilon_k).
\]
Direct multiplication gives
\[
 M^{-1}=\diag(r_1^{-1},\ldots,r_k^{-1})+\frac1{r_0}\mathbf1\mathbf1^*.
\]
The generalized eigenvalues of $(D,M)$ are those of
$D^{1/2}M^{-1}D^{1/2}$, which is $S(\epsilon)$.
This proves Theorem~\ref{thm:main}.

\subsection{Consequences for separated gaps}
The determinant lemma and the trace give
\begin{equation}\label{eq:det-trace}
 \prod_{j=1}^k\nu_j=\frac r{r_0}\prod_{j=1}^k\epsilon_j,
 \qquad
 \sum_{j=1}^k\nu_j=\sum_{j=1}^k\left(1+\frac{r_j}{r_0}\right)\epsilon_j.
\end{equation}
Also, by Cauchy--Schwarz,
\[
 \diag(\epsilon)\le S(\epsilon)
       \le\frac r{r_0}\diag(\epsilon).
\]
If $\epsilon_{(1)}\le\cdots\le\epsilon_{(k)}$ are the ordered
surpluses, then
\begin{equation}\label{eq:gap-bounds}
 \epsilon_{(j)}\le\nu_j\le\frac r{r_0}\epsilon_{(j)}.
\end{equation}
Thus distinct orders of vanishing in the slope surpluses give the
same ordered vanishing orders in the small HE spectrum. The matrix
$S$ determines the constants as well, including repeated scales.

For example, with two leaves and $\epsilon_1/\epsilon_2\to0$,
the larger eigenvalue has leading term
$(1+r_2/r_0)\epsilon_2$. The determinant identity then gives
\[
 \nu_1\sim\frac{r_0+r_1+r_2}{r_0+r_2}\epsilon_1.
\]
The larger edge effectively adds the second leaf's rank to the core
mass before the smaller eigenvalue is formed. This is the elementary
two-edge case of the graph contraction principle in \cite{Flow}.

\section{The second-order term for two factors}
\label{sec:two-factor}

Consider a nonsplit extension $0\to F\to E\to G\to0$ with
$F,G$ nonisomorphic and stable of equal slope at $\omega_0$.
Write $a=\rk F$, $b=\rk G$ and $r=a+b$. In the stable side put
\[
 \delta_s=\kappa_{\omega_s}(G)-\kappa_{\omega_s}(F)>0,
 \qquad N_s=V_s\frac{ab}{r}.
\]
Theorem~\ref{thm:main} gives $\lambda_1\sim\delta_s$.
We now identify the next coefficient for the actual HE metrics.

\subsection{An exact identity at an HE metric}
Use the orthogonal smooth splitting determined by any HE metric on
$E$. Its Dolbeault operator has the form
\[
 \db_E=\begin{pmatrix}\db_F&\beta\\0&\db_G\end{pmatrix}.
\]
The diagonal metrics here are induced metrics; they need not be HE.
Let $P_F$ be the orthogonal projection onto $F$, and set
$q=P_F-(a/r)I$. Then $\|q\|^2=N_s$. The subbundle degree formula
gives
\begin{equation}\label{eq:exact-rayleigh}
 \|\beta\|^2=2\pi a\bigl(\mu(E)-\mu(F)\bigr)
             =N_s\delta_s,
 \qquad \frac{\langle q,\Delta q\rangle}{\|q\|^2}=\delta_s,
\end{equation}
because $\db_{\End}q=-\beta$. This Rayleigh upper bound is a
consequence of the classical degree identity.

The off-diagonal HE equation and the K\"ahler identity give
$\db_{\Hom}^*\beta=0$ for the induced connections. In a unitary
coframe write $\beta=\sum_\nu\beta_\nu\bar\theta^\nu$ and put
\begin{equation}\label{eq:C-beta}
 C_s(\beta)=\diag\left(\sum_\nu\beta_\nu\beta_\nu^*,
                       -\sum_\nu\beta_\nu^*\beta_\nu\right).
\end{equation}
Block multiplication in $\db_{\End}^*\db_{\End}q$ gives
\begin{equation}\label{eq:Delta-q}
 \Delta q=C_s(\beta).
\end{equation}
The term linear in $\beta$ is $-\db_{\Hom}^*\beta$ and vanishes;
the remaining quadratic term is \eqref{eq:C-beta}.

Let $Q$ project onto $q^\perp$ in complex trace-free endomorphisms,
and write $z=C_s(\beta)-\delta_sq$. Suppose the compression
$A=Q\Delta Q$ has spectrum at least $\gamma>\delta_s$.
For the unique small eigenvector, take its $q$ component to be one.
The complementary equation gives $v=-(A-\lambda)^{-1}z$.
Pairing the remaining equation with $q$ yields
\begin{equation}\label{eq:exact-schur-two}
 \lambda=\delta_s-\frac{\langle z,(A-\lambda)^{-1}z\rangle}{N_s}.
\end{equation}
In particular
\[
 \delta_s-\frac{\|z\|^2}{N_s(\gamma-\delta_s)}
       \le\lambda\le\delta_s.
\]

\subsection{The normalized extension limit}
Let $\beta_s$ be the harmonic representative for the split factor
HE metrics, and $B_s=\|\beta_s\|^2$. The scalar equation in
Section~\ref{sec:construction}, now paired with $q$, reads
\[
 -N_s\delta_s+B_st^2+O(t^4)=0.
\]
The coefficient $B_s$ is exact: the quadratic metric correction is
in the image of $L_s$ and pairs trivially with its parallel kernel
vector $q$. Parity makes the scalar equation even, and its derivative
in $w=t^2$ is $B_s+O(w)>0$. Consequently
\begin{equation}\label{eq:two-scaling}
 t_s^2=\frac{N_s}{B_s}\delta_s+O(\delta_s^2).
\end{equation}
These estimates are uniform in $s$, including flat approaches to
$\delta=0$.

We need an orthogonal identification that preserves the subbundle.
The involution $J=\diag(-I_F,I_G)$ sends $t$ to $-t$; uniqueness
gives $u(s,-t)=Ju(s,t)J$. Together with $u_t(s,0)=0$, this makes
the diagonal entries even of order $O(t^2)$ and the off-diagonal
entries odd of order $O(t^3)$. Write
\[
 e^u=\begin{pmatrix}A&B\\B^*&D\end{pmatrix},\qquad
 T=\begin{pmatrix}I&-A^{-1}B\\0&I\end{pmatrix}.
\]
Then
\[
 T^*e^uT=\diag(A,D-B^*A^{-1}B).
\]
Thus
$S=T\diag(A^{-1/2},(D-B^*A^{-1}B)^{-1/2})$ is an isometry
from the split metric to the corrected metric and preserves $F$.
Its diagonal correction is $O(t^2)$ and its upper entry is $O(t^3)$
in every $C^m$ norm. Including the derivative of $S$ in
$S^{-1}\db_{s,t}S$, the upper Dolbeault entry is therefore
\[
 \beta_{\mathrm{HE},s,t}=t\beta_s+O_{C^m}(t^3).
\]
The estimates at one higher spatial derivative justify this formula
in $C^m$. The diagonal connections converge smoothly to the split
wall connections. With \eqref{eq:two-scaling} we obtain
\begin{equation}\label{eq:normalized-beta}
 \frac{\beta_{\mathrm{HE},s}}{\sqrt{\delta_s}}
     \longrightarrow\beta_*:=\sqrt{\frac{N_0}{B_0}}\,\beta_0
 \quad\hbox{in every }C^m,\qquad \|\beta_*\|^2=N_0.
\end{equation}

This convergence takes place in the rescaled presentation of the
extension. On the original fixed $E$, the isomorphism to the scaled
extension multiplies $F$ by $t$; the pulled-back metric on $F$ is
therefore $t^2h_F$ to leading order. Smooth convergence to the split
metric is not asserted in that original presentation.

\subsection{The Green coefficient and eigensection}
Let $\Delta_0$ be the endomorphism Dolbeault Laplacian at the split
wall metric, and $G_0$ its inverse on the complement of $q_0$ in
trace-free endomorphisms. Define
\begin{equation}\label{eq:D-coefficient}
 z_0=C_0(\beta_*)-q_0,\qquad
 D=\frac{\langle z_0,G_0z_0\rangle}{N_0}.
\end{equation}
The vector $z_0$ is orthogonal to $q_0$, since
$\langle q_0,C_0(\beta_*)\rangle=\|\beta_*\|^2=N_0$.

\begin{theorem}\label{thm:second-order}
As $s\to0$ through $\delta_s>0$,
\begin{equation}\label{eq:second-order}
 \lambda_1(s)=\delta_s-D\delta_s^2+o(\delta_s^2).
\end{equation}
Here $D\ge0$, with equality precisely when $C_0(\beta_*)=q_0$
pointwise. After smooth unitary identifications with the wall spaces,
the normalized first eigensection, with positive $q_s$ component, is
\begin{equation}\label{eq:eigensection}
 \psi_s=\frac{q_s}{\sqrt{N_s}}
       -\delta_s\frac{G_0z_0}{\sqrt{N_0}}+o_{C^m}(\delta_s)
 \quad\hbox{for every }m.
\end{equation}
All other trace-free eigenvalues remain uniformly positive.
\end{theorem}

\begin{proof}
Use the orthogonal presentation above and an $L^2$ isometry $U_s$
incorporating the varying volume density. We use $q_s,Q_s,\Delta_s$
for the transported endomorphism objects. Nonlinear expressions must
be formed before this transport: in particular the transported
$C_s(\beta)$ means $U_sC_s(\beta)$, not $C_0$ applied to a
transported form. Density multiplication does not commute with the
quadratic operation $C$.

All differential coefficients converge smoothly to those of $\Delta_0$.
The latter has trace-free kernel $\CC q_0$. Its complementary gap
persists for $A_s=Q_s\Delta_sQ_s$. Indeed, a sequence of normalized
complementary sections with energy tending to zero would have bounded
$H^1$ norm by the first-order elliptic estimate, a strongly convergent
$L^2$ subsequence, and a nonzero limit in
$\ker\Delta_0\cap q_0^\perp$, a contradiction.
Min--max and \eqref{eq:exact-rayleigh} now give exactly one eigenvalue
$\lambda_s<\gamma/2$ for small $s$, and positivity follows from
simplicity of $E$.

Equations \eqref{eq:Delta-q} and \eqref{eq:normalized-beta} imply
\[
 z_s=Q_s\Delta_sq_s=U_sC_s(\beta_{\mathrm{HE},s})-\delta_sq_s,
 \qquad z_s/\delta_s\longrightarrow z_0
\]
in every Sobolev norm. After smooth fixed-complement identifications,
uniform elliptic estimates give
\[
 \|v\|_{H^{j+2}}\le C_j\bigl(\|(A_s-\lambda_s)v\|_{H^j}
                                      +\|v\|_{L^2}\bigr).
\]
The complementary gap absorbs the last term. The projections have
finite rank with smooth kernels and are uniformly bounded on every
Sobolev space. Hence
\[
 \|(A_s-\lambda_s)^{-1}\|_{H^j\to H^{j+2}}\le C_j.
\]
The resolvent identity and coefficient convergence from $H^{j+2}$
to $H^j$ show that these inverses converge to $G_0$ in the stated
operator norm. Substituting into \eqref{eq:exact-schur-two} proves
\eqref{eq:second-order}.

Before normalization the small eigensection is $q_s+v_s$, where
\[
 v_s=-(A_s-\lambda_s)^{-1}z_s
     =-\delta_sG_0z_0+o_{H^{j+2}}(\delta_s).
\]
Its squared norm is $N_s+\|v_s\|^2$. Normalization therefore changes
the $q_s$ coefficient only by $O(\delta_s^2)$. Sobolev embedding,
with arbitrary $j$, gives \eqref{eq:eigensection}; the phase is fixed
by $\langle q_s,\psi_s\rangle_s>0$. Positivity of $G_0$ proves
the assertion about $D$ and its zero set.
\end{proof}

The coefficient is unchanged by nonzero scalar rescaling of the
extension representative, after the induced unitary identifications.
It depends on the wall geometry and the projective extension class.
Ranks and Chern classes alone do not determine it. To compare with
\cite{Flow}, write $c=B_0/N_0$ and
$z=C_0(\beta_0)-cq_0$. Its fixed harmonic-family coefficient is
$d=\langle z,G_0z\rangle/N_0$. Since $z_0=z/c$, we have exactly
$D=d/c^2$. Theorem~\ref{thm:second-order} identifies this coefficient
for the actual HE metrics and the exact slope-gap parameter.

\section{Toric tangent bundles}
\label{sec:toric}

The geometric inputs in this section are proved in the companion
\cite{Chambers}. We state them explicitly to distinguish the geometry
needed for Theorem~\ref{thm:main} from the analytic conclusion.

For $k\ge1$ let
\[
 X_k=\PP_{(\PP^1)^{2k-1}}
       \bigl(\mathcal O\oplus\mathcal O(1^k,-1^{k-1})\bigr),
\]
using quotient projectivization. The twisting line bundle has degree
$1$ on the first $k$ base factors and degree $-1$ on the remaining
$k-1$. Let $\pi:X_k\to(\PP^1)^k$ be the projection onto the
first $k$ base factors, with components $\pi_j$.
This is a smooth toric Fano variety of dimension and Picard number
$2k$. Put $D=-K_{X_k}$ and $h_j=c_1(\pi_j^*\mathcal O_{\PP^1}(1))$.
The tangent bundle fits into
\begin{equation}\label{eq:toric-extension}
 0\longrightarrow F_0:=T_{X_k/(\PP^1)^k}
 \longrightarrow T_{X_k}
 \longrightarrow\bigoplus_{j=1}^k\pi_j^*\mathcal O_{\PP^1}(2)
 \longrightarrow0.
\end{equation}
At $D$, $F_0$ is stable of rank $k$, the quotient line bundles are
stable of the same slope, and all factors are pairwise nonisomorphic.
Every component of the extension class is nonzero. The functions
\[
 z_j(H)=\left(2h_j-\frac{D}{2k}\right)H^{2k-1}
\]
extend to independent real analytic local coordinates at $D$.
The local stable locus is $z_j>0$ for every $j$.

Choose smooth K\"ahler representatives of the nearby ample classes.
The slopes in \cite{Chambers} are intersection numbers without a
factorial denominator. In the present normalization,
\begin{equation}\label{eq:toric-epsilon}
 \epsilon_j(H)=\frac{2\pi z_j(H)}{(2k-1)!V_H}
              =\frac{4\pi k\,z_j(H)}{H^{2k}}.
\end{equation}
At the wall the differential of this map is a common positive
multiple of $dz_j$, so the $\epsilon_j$ are also independent local
coordinates after completion by $k$ further coordinates.

\begin{corollary}\label{cor:toric}
For the HE tangent bundle of $X_k$ in its nearby stable chamber,
there are exactly $k$ small trace-free Dolbeault eigenvalues. To
relative error $O(\sum_j\epsilon_j)$ they are the eigenvalues of
\[
 \diag(\epsilon_1,\ldots,\epsilon_k)+\frac1k vv^*,
 \qquad v_j=\sqrt{\epsilon_j}.
\]
Along a path with every $\epsilon_j=\epsilon$, the leading spectrum
is $\epsilon$ with multiplicity $k-1$ and $2\epsilon$ once.
Positive paths with arbitrary distinct orders of vanishing of the
$\epsilon_j$ are realized by genuine nearby polarizations.
\end{corollary}

\begin{proof}
Apply Theorem~\ref{thm:main} with $r_0=k$ and $r_j=1$.
For equal surpluses the rank-one perturbation has eigenvalue
$\epsilon$ on the complement of $(1,\ldots,1)$ and $2\epsilon$
on its span. Independent local coordinates realize the indicated
positive paths while staying in the ample cone.
\end{proof}

\begin{example}
For $X_2$, write the positive surpluses as $\epsilon_1,\epsilon_2$.
The two comparison eigenvalues are
\[
 \nu_\pm=\frac34(\epsilon_1+\epsilon_2)
    \pm\frac14\sqrt{9(\epsilon_1-\epsilon_2)^2+4\epsilon_1\epsilon_2}.
\]
If $\epsilon_1/\epsilon_2\to0$, the actual HE eigenvalues satisfy
$\lambda_1\sim\frac43\epsilon_1$ and
$\lambda_2\sim\frac32\epsilon_2$. The relative theorem applies
equally when $\epsilon_1=e^{-1/\epsilon_2^2}$.
\end{example}

\begin{example}
For $X_3$ choose $\epsilon_1=s^6$, $\epsilon_2=s^4$ and
$\epsilon_3=s^2$, $s>0$. The characteristic polynomial of the
comparison matrix is
\[
 \lambda^3-\frac43(\epsilon_1+\epsilon_2+\epsilon_3)\lambda^2
 +\frac53(\epsilon_1\epsilon_2+\epsilon_1\epsilon_3+
                         \epsilon_2\epsilon_3)\lambda
 -2\epsilon_1\epsilon_2\epsilon_3.
\]
Its largest root is asymptotic to $\frac43s^2$. The second elementary
symmetric function then gives the middle root $\frac54s^4$, and
the determinant gives the smallest root $\frac65s^6$.
Consequently the actual tangent-bundle spectrum satisfies
\[
 \lambda_1\sim\tfrac65s^6,\qquad
 \lambda_2\sim\tfrac54s^4,\qquad
 \lambda_3\sim\tfrac43s^2.
\]
These are changes of polarization on one fixed six-dimensional variety.
\end{example}

Theorem~4.1 and equation~(16) of \cite{Chambers} supply further examples
from the strictly semistable smooth toric Fano census in dimensions
three through six, overlapping this family at $X_2$ and $X_3$.
All 31 indecomposable cases give star extensions of distinct locally
free stable factors, with nonzero components. Twenty-six have two
factors and hence also
satisfy Theorem~\ref{thm:second-order}. The five cases with more
factors have ranks
\[
 (2,1,1),\ (3,1,1),\ (4,1,1),\ (4,1,1),\ (3,1,1,1).
\]
In dimension $n$, for a quotient of rank $r_j$ with first Chern class
$C_j$, the general conversion is
\[
 \epsilon_j(H)=\frac{2\pi}{r_j(n-1)!V_H}
       \left(C_j-\frac{r_j}{n}c_1(T_X)\right)H^{n-1}.
\]
Computing the two-factor
coefficient $D$ would additionally require the wall Green operator
and normalized harmonic extension; the toric degree data do not
provide those analytic quantities.

\section{Scope and further questions}

The star structure is responsible for the disappearance of harmonic
norms from the leading spectrum. It gives one positive scaling
variable for each independent scalar equation, and a split summand
on every coordinate face. For general two-step extension graphs
with cycles, the number of edge weights can exceed the number of
independent moment equations; the ranks and slope surpluses then
need not determine all leading weights. Repeated stable factors
also replace block scalars by noncommutative endomorphism algebras.
Neither setting is included in Theorem~\ref{thm:main}.

The estimates concern equilibrium metrics for varying K\"ahler forms.
Relating them to the long-time HYM flow as the polarization also
approaches a wall requires a uniform two-parameter analysis. The
inverse-time behavior on the wall and the positive equilibrium gap
inside the chamber do not by themselves justify interchanging those
limits.


\begin{thebibliography}{Kob87}
\bibitem[CT24]{CT} A. Clarke and C. Tipler,
\emph{Semi-stability and local wall-crossing for Hermitian Yang--Mills
connections}, J. Symplectic Geom. \textbf{22} (2024), no.~3, 673--694;
arXiv:2304.05245.
\bibitem[Del25]{Delloque} R. Delloque,
\emph{Continuity of HYM connections with respect to metric variations},
J. Lond. Math. Soc. (2) \textbf{112} (2025), no.~1, e70219;
arXiv:2403.16814.
\bibitem[Kob87]{Kobayashi} S. Kobayashi,
\emph{Differential Geometry of Complex Vector Bundles}, Publications
of the Mathematical Society of Japan, vol.~15, Princeton University
Press and Iwanami Shoten, 1987.
\bibitem[ST24]{ST} L. M. Sektnan and C. Tipler,
\emph{Hermitian Yang--Mills connections on pullback bundles},
Calc. Var. Partial Differential Equations \textbf{63} (2024),
Paper No.~13; arXiv:2006.06453v2.
\bibitem[UY86]{UY} K. Uhlenbeck and S.-T. Yau,
\emph{On the existence of Hermitian--Yang--Mills connections in stable
vector bundles}, Comm. Pure Appl. Math. \textbf{39} (1986),
suppl., S257--S293.
\bibitem[Wue26a]{Flow} B. J. Wuebben,
\emph{Small Dolbeault eigenvalues along Hermitian Yang--Mills flow},
preprint, 2026, \url{https://doi.org/10.5281/zenodo.23643223}.
\bibitem[Wue26b]{Chambers} B. J. Wuebben,
\emph{Local stability chambers for toric tangent bundles},
preprint, 2026, \url{https://doi.org/10.5281/zenodo.23681809}.
\end{thebibliography}
\end{document}
